% TOP_PROBLEM: {"id":30007688,"problem_number":"LOCAL-30007688","title":"Privacy valuation","category_id":21,"category":{"id":21,"name":"economics","display_name":"Economics","description":"Open economic theory statements, published research agendas, and proposed scoped specifications; question provenance and historical priority are qualified per record.","slug":"economics","order_index":21,"created_at":"2026-10-08T00:00:00Z"},"status":"provenance_pending","external_url":null,"legacy_ids":[],"published":false,"statement":"Identify or bound informed compensation required for auditable data sharing when choices also reflect misunderstanding and uncertain future harms.","economics":{"original_id":177,"field":"Digital, platform and data economics","kind":"identification","formalization_basis":"proposed_specification","source_status":"not_verified","priority_status":"not_established","priority_note":"No exact open-problem passage was verified. The supporting publication does not establish priority or unresolved status.","earliest_verified_reference":null,"current_reference":{"authors":["Avi Goldfarb","Verina F. Que"],"title":"The Economics of Digital Privacy","year":2023,"url":"https://www.nber.org/papers/w30943","doi":"10.3386/w30943","locator":"Primary publication abstract; full open-question paragraph not retrieved","evidence_summary":"Reviews privacy benefits, costs and data-flow externalities; no exact open statement supporting this specification verified."},"formalization_note":"The cited primary work supports the subject or supplies solutions in particular settings, but no exact open-question passage for this specification was verified. The mathematical target and restrictions are proposed here, with no canonical-conjecture or priority claim. Independent math review tightened feasibility, existence, or estimand assumptions where required."},"statusQualification":"Provenance pending: an explicit published statement of this exact open question was not verified. This is a catalogue-authored candidate specification, not a certified open conjecture. The cited primary work supports the subject or supplies solutions in particular settings, but no exact open-question passage for this specification was verified. The mathematical target and restrictions are proposed here, with no canonical-conjecture or priority claim. Independent math review tightened feasibility, existence, or estimand assumptions where required.","statusReviewedAt":"2026-10-08"}
% FIRST_POSTED: 2026-10-08
% ENTRY_KIND: statement_only
% !TeX program = lualatex
\documentclass[11pt]{article}
\usepackage{fontspec}
\usepackage[a4paper,margin=25mm]{geometry}
\usepackage{amsmath,amssymb,mathtools}
\usepackage{xurl}
\usepackage[hidelinks]{hyperref}
\newcommand{\sref}[2]{\hyperlink{source:#1:#2}{[S#2]}}
\setlength{\emergencystretch}{3em}
\begin{document}
% problemId:30007688
\section*{Privacy valuation}\label{30007688}
\subsection{Definitions and mathematical statement}
Consider adult users of an online service in a fixed country and enrollment period. A data-sharing contract \(d\in\{0,1\}\) either limits collection to essential service data or permits a specified additional use; the contract's actual data flows are auditable. Let \(u_i(c,d,\ell)\) be user \(i\)'s utility from baseline consumption \(c>0\), sharing regime \(d\), and future loss \(\ell\), with utility strictly increasing in \(c\). The additional use induces an uncertain loss distribution \(F_i\) over a fixed five-year horizon, including discrimination and unwanted disclosure. Define compensating variation \(CV_i\) by \(E_{F_i}[u_i(c+CV_i,1,\ell)]=u_i(c,0,0)\), with expectation over the user's stated informed loss beliefs. Require finite expected utility and \(c+CV_i>0\); report nonexistence if no finite compensation attains the comparison utility. The target is the population distribution of \(CV_i\), alongside sensitivity to alternative credible loss distributions. Incentivized contract choices, comprehension checks and belief elicitation provide distinct observations; default effects cannot automatically be interpreted as valuations. Determine which additional information treatments, repeated choices or structural restrictions distinguish preferences from misunderstanding and friction. When distant harms cannot be learned from the observation window, bound valuations under explicit probability and loss limits. This is a proposed identification exercise rather than a verified universal unresolved privacy-valuation theorem.

\par\textbf{Formulation and scope.} The cited primary work supports the subject or supplies solutions in particular settings, but no exact open-question passage for this specification was verified. The mathematical target and restrictions are proposed here, with no canonical-conjecture or priority claim. Independent math review tightened feasibility, existence, or estimand assumptions where required.

\par\textbf{Open-question provenance.} An explicit published open-question statement was not verified. Sources below are supporting literature and must not be treated as a first listing.

\par\textbf{Historical priority.} No exact open-problem passage was verified. The supporting publication does not establish priority or unresolved status.

\subsection{Short English statement}
Identify or bound informed compensation required for auditable data sharing when choices also reflect misunderstanding and uncertain future harms.

\subsection{Sources}
\begin{itemize}\raggedright
\item[\textnormal{[S1]}]\hypertarget{source:30007688:1}{} Avi Goldfarb, Verina F. Que. 2023. The Economics of Digital Privacy. Primary publication abstract; full open-question paragraph not retrieved.
\url{https://www.nber.org/papers/w30943}.
DOI: 10.3386/w30943.
{\small\textit{Supporting or current literature; see evidence qualification. Reviews privacy benefits, costs and data-flow externalities; no exact open statement supporting this specification verified.}}
\end{itemize}
\end{document}
